% Users, roles and the PHI permission; managing accounts; signing in (in brief: the methods are
% chap:sign-in); Privacy Mode for administrators: who is in it, the switch, its settings and the
% display of dates; the audit log and its seals. How the privacy gate de-identifies answers is
% chap:protecting; the security review chap:security.
% Sources: docs/administration.md, docs/sign-in.md, docs/protecting-patient-data.md,
% docs/security.md (5, 6, 12); the code is the authority: medorama/src/users.rs, server.rs
% (in_privacy, set_privacy, user_add), signin.rs, deid.rs (FIELDS, Settings), audit.rs, display.rs.
\chapter{Users and privacy}
\label{chap:users-privacy}

This chapter is for administrators. It describes the accounts of a catalog and what each may do,
how users sign in, how Privacy Mode is set up, and the audit log that records who did what.
Every person who uses Medorama has an account, with a \emph{role} that says what they may do and,
separately, the \emph{PHI permission} that says whether they may see identifiers. A session
without that permission, or whose user has switched it on, is in \textbf{Privacy Mode}: the
server then sends no identifier at all. How the privacy gate de-identifies each answer and file is
explained in \cref{chap:protecting}, and the standards Medorama follows and its limits are
reviewed in \cref{chap:security}.

\section{Users and roles}
\label{sec:users-privacy-roles}

Each account has one of three roles (\cref{tab:users-privacy-roles}).

\begin{xltabular}{\linewidth}{@{}P{2.2cm}L@{}}
\caption{The roles. Each role may do everything the one above it may.}
\label{tab:users-privacy-roles}\\
\toprule
\tabhead Role & \tabhead May \\
\midrule
\endfirsthead
\toprule
\tabhead Role & \tabhead May \\
\midrule
\endhead
\bottomrule
\endfoot
\code{viewer} & Search and look at patients, encounters, recordings, databases and results; keep
  a history; make, keep and share cohorts \\
\code{analyst} & Also run and cancel actions (indexing, plug-ins) \\
\code{admin} & Also import and remove data, connect databases, review identities, manage users
  and settings (Privacy Mode, dates and times), and read the audit log \\
\end{xltabular}

Every check is made on the server. The page hides what a user may not do, but each request is
checked again. Every signed-in user may see every patient of the catalog, as Privacy Mode and
their permission allow: there are no per-patient access lists, and separate studies are served by
separate catalogs. Cohorts are private to their owner unless they are shared
(\cref{chap:cohorts}).

\subsection{The PHI permission}
\label{sec:users-privacy-phi}

The PHI permission is set for each user, on top of the role. It lets the user leave Privacy Mode,
and then see and search MRNs, names, dates of birth, account and encounter numbers and real
dates, and read the identifiers that recordings' files state. Without it, every session of the
user is in Privacy Mode, an administrator's too. With it, a session starts in Privacy Mode, except
an administrator's, which starts outside it (\cref{tab:protecting-who}). The first administrator,
created by \cmd{medorama init}, has the permission; any other user has it only when an
administrator gives it.

\subsection{Managing users}
\label{sec:users-privacy-managing}

Users are managed in \ui{Admin} $\rightarrow$ \ui{Users}, or on the server with
\cmd{medorama user} (\cref{chap:cli}). The following commands add two users, reset a forgotten
password, disable an account and list the accounts:
\begin{lstlisting}[style=shell]
medorama user add alice --role analyst            # asks for her password
medorama user add bob --role viewer --phi         # may see identifiers
medorama user set alice --password                # a forgotten password
medorama user set bob --disabled true             # cannot sign in; his sessions end
medorama user list
\end{lstlisting}
When users sign in through the institution's sign-in (\code{sign\_in = "proxy"}), accounts need
no password: \cmd{medorama user add NAME -{}-role R -{}-no-password} gives the account a random
password that nobody knows, and \ui{Admin} $\rightarrow$ \ui{Users} hides the password field and
the \ui{Reset password} button (\cref{sec:sign-in-accounts}).
\begin{itemize}
\item \textbf{User names} consist of letters, digits and \code{. \_ - @}, at most 64 characters.
  With proxy sign-in, they are the names the proxy sends.
\item \textbf{Passwords} have at least 10 characters and are stored as Argon2id hashes
  (\cref{sec:sign-in-password}).
\item \textbf{At least one enabled administrator.} The catalog refuses a change or a removal that
  would leave it without one (\code{the catalog needs at least one enabled admin}).
\item \textbf{Every change is audited} (\code{user\_added}, \code{user\_changed},
  \code{user\_deleted}, \code{password\_changed}).
\end{itemize}

\section{Signing in}
\label{sec:users-privacy-signin}

A catalog signs users in with one of two methods, chosen by the setting \code{sign\_in}.
With \code{password}, the default, users sign in on Medorama's own form, with throttled attempts
and sessions that end after \code{session\_hours} (default 12) or when unused for
\code{session\_idle\_minutes} (default 60). With \code{proxy}, the institution's authenticating
reverse proxy signs users in, with single sign-on and multi-factor authentication where the
institution has them, and names each user to Medorama in a header. Both methods, their settings,
the proxy's security requirements and example configurations are described in
\cref{chap:sign-in}.
\begin{itemize}
\item \textbf{Signing out everywhere.} In password mode, a user ends all their sessions, on every
  computer and browser, with the user menu's \ui{Sign out everywhere} (audited as
  \code{signed\_out\_everywhere}). An administrator ends a user's sessions with \ui{Admin}
  $\rightarrow$ \ui{Users} $\rightarrow$ \ui{Sign out} (audited as \code{user\_changed}). A new
  password or a disabled account also ends them, and a disabled user is refused at the next
  request.
\item \textbf{The command line} acts as the catalog's owner, with administrator rights and the
  PHI permission, and is recorded in the audit log as \code{local:<user>}. Whoever can run it
  under the catalog's account can read the database directly anyway.
\end{itemize}

\section{Privacy Mode}
\label{sec:users-privacy-mode}

The catalog keeps every identifier as imported, and the server decides, at the moment of
sending, what leaves. In Privacy Mode, no identifier reaches the browser: the page never receives
what it must not show, and it does no masking of its own. Who is in Privacy Mode is summarized in
\cref{tab:protecting-who}, and what the privacy gate replaces, moves and withholds is explained in
\cref{chap:protecting}.

The \ui{Privacy Mode} switch in the page's header sends \code{POST privacy}, which is audited
(\code{privacy\_on}, \code{privacy\_off}). The choice is kept by the server for the session,
never in the browser: a new sign-in starts in the user's default mode
(\cref{sec:sign-in-password,sec:sign-in-proxy}). The page then reloads, so that nothing
identified stays on the screen. A user without the permission who asks
to leave Privacy Mode is refused (HTTP 403). A session in Privacy Mode is answered as if its user
had no PHI permission, and its answers pass the privacy gate.

\subsection{Privacy Mode's settings}
\label{sec:users-privacy-settings}

\ui{Admin} $\rightarrow$ \ui{Privacy Mode} chooses what Privacy Mode withholds and how far it
moves dates (\cref{tab:users-privacy-fields}). The defaults follow HIPAA's Safe Harbor list where
the catalog holds the data, except the text of notes and reports, which is shown with a warning.
Every change is audited (\code{deid\_settings}). Only administrators may read these settings
(\code{GET deid}), the date offset among them; other users learn from \code{GET config} what is
withheld, never the offset.

\begin{xltabular}{\linewidth}{@{}P{3.2cm}L P{1.9cm}@{}}
\caption{What Privacy Mode withholds: the choices of \ui{Admin} $\rightarrow$ \ui{Privacy Mode}.}
\label{tab:users-privacy-fields}\\
\toprule
\tabhead Choice & \tabhead Covers & \tabhead Default \\
\midrule
\endfirsthead
\toprule
\tabhead Choice & \tabhead Covers & \tabhead Default \\
\midrule
\endhead
\bottomrule
\endfoot
\code{names} & Patients' names & Withheld \\
\code{birth\_date} & Dates of birth & Withheld \\
\code{dates} & Real dates and times, moved back by the offset & Moved \\
\code{age\_over\_89} & Ages above 89, shown as \code{90+} & Withheld \\
\code{mrn} & Medical record numbers & Withheld \\
\code{account\_numbers} & Account and encounter numbers (CSN) & Withheld \\
\code{notes} & The text of notes and reports (\ui{The text of notes and reports}), and the
  annotations of EDF+ recordings in VitalStream. When it is withheld, notes read
  \code{[Withheld]} and annotations \code{[Annotation]}. Otherwise the text is shown with the
  patient's own names and date of birth replaced, and notes are marked with a warning, because
  free text may still name the patient or others & Shown, marked \\
\end{xltabular}

Some protections do not depend on these choices and always apply: imported patient and encounter
IDs become pseudonyms, file names become tokens, staff names become \code{[Staff Name]},
unrecognized columns are withheld, and files that the gate cannot de-identify are refused
(\cref{sec:protecting-settings}).

\paragraph{Dates.} Every date and time moves back by \textbf{the catalog's offset}, written in
years, months, days and hours (each 0 or more) and counted on the calendar from 2000-01-01; the
default is 36\,570 days, about 100 years. Each patient's dates also move back by \textbf{the
patient's own days}, a whole number from 0 to \code{spread\_days} (default 365, at least 30, at
most 3650), drawn at random once, when the patient's pseudonym is drawn, and kept with it. A
\code{spread\_days} under 30 saved before that minimum is kept, since raising it moves every
patient's dates, and \cmd{medorama serve} warns about it at start. The intervals between one patient's events are therefore kept exactly
(\cref{sec:protecting-dates}).

\paragraph{Dates and times on the page.} \ui{Admin} $\rightarrow$ \ui{Dates and times} chooses how
every page writes dates and times: \code{2024-01-31} (the default), \code{01/31/2024},
\code{31/01/2024} or \code{31.01.2024}, with a 24-hour (the default) or 12-hour clock. Exports and
the API keep ISO 8601. Every change is audited (\code{display\_changed}).

\begin{note}
Record the offset and the spread under your governance, and do not disclose the offset. Anyone
who learns the real date of one of a patient's events can re-date every event of that patient,
and every other patient's to within \code{spread\_days}. Date shifting is pseudonymization, not
anonymization (\cref{sec:protecting-limits}).
\end{note}

\section{The audit log}
\label{sec:users-privacy-audit}

The audit log is append-only and readable by administrators (the \ui{Admin} tab,
\cmd{medorama audit}, \code{GET audit}). Each entry records the time, the user, the event, its
target and a detail (\cref{tab:users-privacy-events}). In Privacy Mode, the audit page shows no
target or detail, because a search's text may name a patient. Medorama never deletes an entry.
An entry is never lost when the database cannot take it at once (busy beyond its wait, or
failing): it is kept, in order, and written before the next entry or by the next seal; with
\code{audit\_export} set it is appended to that file at once (a line with
\code{"pending": true}), and the server's output says so.

\begin{xltabular}{\linewidth}{@{}P{3.0cm}L@{}}
\caption{The events of the audit log.}
\label{tab:users-privacy-events}\\
\toprule
\tabhead Area & \tabhead Events \\
\midrule
\endfirsthead
\toprule
\tabhead Area & \tabhead Events \\
\midrule
\endhead
\bottomrule
\endfoot
Sign-in & \code{login} (detail \code{proxy} with proxy sign-in, one per visit), \code{login\_failed}, \code{password\_changed},
  \code{signed\_out\_everywhere} (a user's own \ui{Sign out everywhere}) \\
Users and settings & \code{user\_added}, \code{user\_changed} (an administrator's \ui{Sign out} of
  another user too, written in its detail), \code{user\_deleted}, \code{deid\_settings},
  \code{display\_changed}, \code{init}; on the command line, \code{addon\_install}
  and \code{addon\_remove} \\
Privacy Mode & \code{privacy\_on}, \code{privacy\_off} \\
Identifiers & \code{phi\_search} (searches of identifiers and of notes' text),
  \code{privacy\_search} (searches, filters and cohort criteria made in Privacy Mode, with their
  text), \code{phi\_view},
  \code{patient\_view}, \code{phi\_export} (an export by a user with the PHI permission outside
  Privacy Mode) \\
Data & \code{import\_table}, \code{import\_recordings}, \code{cancel\_import},
  \code{remove\_import}, \code{remove\_encounter}, \code{remove\_recording},
  \code{remove\_matching}, \code{merge\_patients}, \code{distinct\_patients},
  \code{move\_encounter}, \code{link\_decision}, \code{apply\_held\_row},
  \code{dismiss\_held\_row}, \code{view\_database}, \code{export\_database}, \code{view} (a
  recording opened), \code{demo\_add}, \code{demo\_remove} \\
Connected databases & \code{connect}, \code{disconnect}, \code{link} (persons linked, held,
  invalid), \code{copy\_database} \\
Cohorts & \code{cohort\_create}, \code{cohort\_change}, \code{cohort\_refresh},
  \code{cohort\_delete} (each with what changed, patients by ID), \code{export} (a cohort's
  members' data) \\
Work & \code{run}, \code{check} \\
\end{xltabular}

\subsection{Seals}
\label{sec:users-privacy-seals}

The log is tamper-evident, as follows.
\begin{itemize}
\item \textbf{Seals.} Every minute, the server seals the entries written since the last seal. A
  seal is the SHA-256 chain of those entries, in order, starting from the previous seal's hash,
  and each entry records the seal that covers it. Entries committed late (in a long transaction) go
  into the next seal.
\item \textbf{Verification.} \cmd{medorama audit -{}-verify} (or \code{GET audit/verify}, for
  administrators) computes every chain again. An entry changed, deleted or added to a sealed batch, or a
  seal changed or deleted, shows as the first seal that no longer holds, and the command fails.
  Entries and seals are numbered by the database without reuse, so verification also counts
  what it numbered against what is left: the last seal deleted with its entries, or entries
  deleted before they were sealed, show as missing.
  \cmd{medorama audit -{}-seal} seals now, and \cmd{medorama audit -{}-seal -{}-verify} seals,
  then checks every seal, the new one too.
\item \textbf{Out of reach.} Whoever controls the database could rewrite entries and seals
  together. With \code{audit\_export} set, each seal and its entries are also appended to that
  file (JSON lines, readable by its owner only, synced to disk): put it on storage the database's
  administrators cannot change (a write-once volume, a log collector). Verification compares the
  database's seals with it (each exported seal must be in the database, with the same hash): the
  end of the log is protected by the export, up to its last seal; entries written since are
  protected by the database's numbering only.
\end{itemize}
\begin{lstlisting}[style=shell]
medorama audit --limit 20          # the latest entries
medorama audit --verify            # whether every seal holds
\end{lstlisting}
